\documentclass[11pt,a4paper]{article}
\usepackage[margin=27mm]{geometry}
\usepackage{amsmath,amssymb,amsthm,fontspec,unicode-math,hyperref,booktabs,longtable}
\usepackage{array,calc,fvextra}
\setmainfont{STIXGeneral}
\setmathfont{STIX Two Math}
\setmonofont{Menlo}
\hypersetup{hidelinks}
\setlength{\emergencystretch}{2em}
\providecommand{\tightlist}{\setlength{\itemsep}{0pt}\setlength{\parskip}{0pt}}
\newtheorem{theorem}{Theorem}
\newtheorem{lemma}[theorem]{Lemma}
\newtheorem{corollary}[theorem]{Corollary}
\newcommand{\gen}[1]{\langle #1\rangle}
\newcommand{\Sym}{\operatorname{Sym}}
\newcommand{\Aut}{\operatorname{Aut}}
\newcommand{\Fix}{\operatorname{Fix}}

\title{Commuting involutions and six-point polyhedra}
\author{Human Compression project}
\date{}
\begin{document}
\maketitle
\begin{abstract}
The five transitive actions of a marked four-group classify pairs of
commuting involutions in symmetric groups. Parity and reversal give the
all-degree formulas of Kiefer and Leemans for symmetric and alternating
groups. In degree six the full automorphism group requires a separate
calculation, which also gives their complete regular-polyhedron
classification. The finite program is exhaustive and is included in full.
\end{abstract}
The equivalence here is the action of the full automorphism group of the
ambient group, with interchange of the two involutions permitted. Except
in degree six, the classical automorphism theorem identifies this with
permutation conjugacy and reversal. The exception is treated explicitly;
it cannot be absorbed into a formula for permutation conjugacy.

The source is A.~Kiefer and D.~Leemans, \emph{On pairs of commuting involutions
in Sym(n) and Alt(n)}, Communications in Algebra \textbf{41} (2013),
4408--4418. We use the previously bound nine-page author edition, not an
unverified identification with every detail of the publisher's typesetting.
The results and their historical attribution remain those of the source.
The shared finite-action and extension theorems are in [UD].

\section{Multiplicities, parity and
centralizers}\label{multiplicities-parity-and-centralizers}

Two distinct nonidentity commuting involutions \(a,c\) generate
\(V=\langle a,c\rangle\cong C_2^2\). Conversely any faithful permutation
action of this marked \(V\) supplies such a pair. Apply UD.Orbits with
orbit stabilizers
\(V,\langle a\rangle,\langle c\rangle,\langle ac\rangle,1\) and
multiplicities \(f,x,y,z,t\), respectively. The complete invariant under
permutation conjugacy is

\[
 N=f+2(x+y+z)+4t,\qquad
 t>0\ \text{or at least two of }x,y,z\text{ are positive}.
 \tag{CPI1}
\]

The second condition is precisely UD.Orbits' trivial-kernel condition.
Interchanging the named generators exchanges \(x,y\) and fixes
\(f,z,t\). Both permutations are even exactly when \(y+z\) and \(x+z\)
are even, equivalently \(x\equiv y\equiv z\pmod2\). Indeed an element
outside an index-two stabilizer flips that two-point orbit; on a regular
four-point orbit it makes two transpositions and is even. Thus only the
three two-point orbit multiplicities contribute to the sign character.
Every tuple occurs as the disjoint union of the indicated coset actions,
so the classification is an existence and uniqueness theorem, not only a
count.

For an involution \(\rho\) with \(k\) transpositions in \(S_N\),
UD.Orbits, equation (5.4) for \(C_2\) gives

\[
 C_{S_N}(\rho)\cong(C_2\wr S_k)\times S_{N-2k}.
 \tag{CPI2}
\]

If \(\rho\in A_N\), its alternating centralizer is the kernel of the
sign on this product. A flip on one two-point orbit is an odd element in
the centralizer. Thus this kernel has index two and order
\(2^{k-1}k!(N-2k)!\). This proves source Propositions 3.1 and 4.1
directly from the shared finite-action theorem. The same odd centralizer
element also shows that a single involution's \(S_N\)-class does not
split in \(A_N\). This statement about individual involutions is not a
replacement for the required pair-orbit equivalence.

\section{The formulas in every nonexceptional
degree}\label{the-formulas-in-every-nonexceptional-degree}

Let \(n=\lfloor N/2\rfloor\). For a pair with \(a\) containing \(k\)
transpositions, let \(l\) be the number of \(c\)-transpositions inside
the support of \(a\). Its multiplicities satisfy the reversible
coordinate change

\[
 \begin{gathered}
 k=y+z+2t,\qquad l=z+2t,\\
 0\le l\le k,\qquad 0\le t\le\lfloor l/2\rfloor,\qquad0\le x\le n-k,\\
 z=l-2t,\qquad y=k-l,\qquad f=N-2(x+k).
 \end{gathered}
 \tag{CPI3}
\]

For \(k\ge1\) the only nonfaithful choices are \(c=1\) (\(x=l=t=0\)) and
\(c=a\) (\(x=t=0,l=k\)). Indeed any other commuting involution distinct
from \(a\) yields a faithful marked \(V\), by (CPI1). Define

\[
 \lambda(k)=\sum_{l=0}^k(\lfloor l/2\rfloor+1)
 =\lfloor(k/2+1)^2\rfloor.
 \tag{CPI4}
\]

The equality follows by pairing consecutive even and odd \(l\); it holds
for both parities of \(k\). Thus the number of ordered conjugacy classes
is

\[
 O(n)=\sum_{k=1}^n\lambda(k)(n-k+1)-2n.
 \tag{CPI5}
\]

This is exactly source Lemma 3.2. Swapping the endpoints fixes a
multiplicity vector precisely when \(x=y\), equivalently the two
permutations have the same transposition count. If those counts differ
the ordered classes are paired by reversal; if equal, reversal fixes the
ordered class itself. This proves source Lemma 3.3, without using its
printed inference from conjugacy of a single second involution to
conjugacy fixing the first.

For a fixed class with common cycle count \(k\), put \(x=y=k-l\). The
lower bound on \(l\) is \(\max(0,2k-n)\). The choice \(c=a\) is
excluded. Set \(\psi(k,n)=0\) for \(k\le\lfloor n/2\rfloor\). For
\(h=2k-n>0\), the omitted initial sum is

\[
 \psi(k,n)=\sum_{l=0}^{h-1}(\lfloor l/2\rfloor+1)
 =\begin{cases}h^2/4+h/2,&n\text{ even},\\(h-1)^2/4+h,&n\text{ odd}.
 \end{cases}
 \tag{CPI6}
\]

This follows from the same even/odd pairing and \(h\equiv n\pmod2\).
Consequently the number fixed at \(k\) is
\(D(k,n)=\lambda(k)-\psi(k,n)-1\). The order-two case of
UD.ExactSupportOrbits counts reversal orbits as half the sum of the
total and fixed counts. It gives

\[
 s_n=-\frac{3n}{2}+\sum_{k=1}^n\lambda(k)\left(\frac{n-k}{2}+1\right)
 -\frac12\sum_{k=\lfloor n/2\rfloor+1}^n\psi(k,n).
 \tag{CPI7}
\]

For even endpoints, \(k\) must be even and \(l+x\) must be even. When
\(k\) is even the numbers of \((l,t)\) with \(l\) even or odd are
respectively

\[
 \gamma(k)=\sum_{j=0}^{k/2}(j+1)=k^2/8+3k/4+1,\qquad
 \delta(k)=\sum_{j=0}^{k/2-1}(j+1)=k^2/8+k/4.
 \tag{CPI8}
\]

Among \(x=0,\ldots,n-k\), exactly \(\lceil(n-k+1)/2\rceil\) are even and
\(\lfloor(n-k+1)/2\rfloor\) odd. The two removed choices are even.
Therefore

\[
 \begin{gathered}
 \mu(n)=-2\lfloor n/2\rfloor+
 \sum_{\substack{1\le k\le n\\k\text{ even}}}
 \left[\gamma(k)\left\lceil\frac{n-k+1}{2}\right\rceil+
 \delta(k)\left\lfloor\frac{n-k+1}{2}\right\rfloor\right],\\
 a_n=\frac12\left(\mu(n)+\sum_{\substack{1\le k\le n\\k\text{ even}}}D(k,n)\right).
 \end{gathered}
 \tag{CPI9}
\]

The first expression proves source Lemma 4.2 and the second is exactly
Theorem 1.2's formula, since its \(\varphi(k,n)\) is \(D(k,n)\). These
formulas count \(S_N\)-conjugacy classes for all \(N\ge4\), including
six. To obtain the source's isomorphism relation we retain the standard
theorem that \(\operatorname{Aut}(S_N)\) is induced by \(S_N\)
conjugation and \(\operatorname{Aut}(A_N)\) by \(S_N\) conjugation for
\(N\ge4\), \(N\ne6\). Thus (CPI7) and (CPI9) prove the nonexceptional
parts of source Theorems 1.1 and 1.2. For \(N=6\) the full automorphism
action is computed below. All nontrivial \(V\)-orbits have even size, so
removing one fixed point in odd degree and adding one in even degree is
a bijection of the multiplicity sets. This proves the even/odd degree
pairing; \(N=6\) and \(N=7\) differ only when the ambient exceptional
automorphisms are allowed.

As a compact alternative encoding of the same classification, let \(v\)
record half-degree. With \(x\le y\) for unordered pairs,

\[
 \begin{aligned}
 \sum_{n\ge0}s_nv^n&=\frac1{(1-v)^3(1-v^2)^2}-\frac{1+v}{(1-v)^2},\\
 \sum_{n\ge0}a_nv^n&=\frac{1+v^3}{(1-v)(1-v^2)^3(1-v^4)}-
 \frac{1+v^2}{(1-v)(1-v^2)}.
 \end{aligned}
 \tag{CPI10}
\]

Here \(s_n,a_n\) continue to denote conjugacy counts, without the
six-point automorphism fusion. To verify every factor, write \(y=x+d\):
\(x,d,z,t\) and unused pairs of fixed points have weights \(2,1,1,2,1\).
The nonfaithful tuples have \(t=0\) and at most one positive coordinate
among \(x,y,z\). Their series is \((1+v)/(1-v)^2\). For even endpoints
write \(x=2x'+e,y=2y'+e,z=2z'+e\) with \(e\in\{0,1\}\). The common odd
shift has weight 3; the other weights are \(4,2,2,2,1\). The nonfaithful
even tuples have \(e=t=0\) and give the second subtraction. This proves
the rational expressions independently of a finite coefficient check.

\section{The exceptional six-point
actions}\label{the-exceptional-six-point-actions}

Apply UD.ExtensionOrbits with \(B=\operatorname{Aut}(G)\), \(E\) the
unordered pairs \(\{a,c\}\) of distinct commuting involutions, and \(Z\)
the pairs together with an involution \(b\). An automorphism acts on the
unordered endpoints and on \(b\); it may interchange \(a,c\). Thus its
orbit on \(Z\) is exactly an isomorphism of the marked triple with
endpoint reversal allowed. This is precisely the equivalence of regular
polyhedra up to isomorphism and duality.

The complete degree-six certificate constructs every involution by
enumerating the elements of \(G\), then every distinct unordered
commuting pair, and its full automorphism orbit. For each orbit it
enumerates every middle involution modulo the endpoint-pair stabilizer.
It tests generation and, for every generating triple, all 64 ordered
pairs of subsets \(I,J\subseteq\{0,1,2\}\) for
\(\langle r_i:i\in I\rangle\cap\langle r_i:i\in J\rangle=\langle r_i:i\in I\cap J\rangle\).
There is no shortcut based solely on group orders, endpoint type, or a
published classification list. The explicitly tested intersection
property and the retained string-C-group/polytope correspondence then
identify the surviving triples with all required polyhedra.

The full required program and the additional formula diagnostic are
printed in the companion program appendix; their exact roles are
specified below. Its exact finite output is:

{
\begin{longtable}[]{@{}lrr@{}}
\toprule\noalign{}
Quantity & \(S_6\) & \(A_6\) \\
\midrule\noalign{}
\endhead
\bottomrule\noalign{}
\endlastfoot
Automorphism-group order & 1440 & 1440 \\
Involutions & 75 & 45 \\
Unordered commuting pairs & 495 & 90 \\
Endpoint orbits & 5 & 1 \\
Middle orbits & 80 & 8 \\
Generating triples & 2 & 0 \\
Intersection-property survivors & 2 & 0 \\
\end{longtable}
}

The two \(S_6\) survivors have unordered adjacent orders \{5,6\} and
\{6,6\}. Exact representatives, separately checked for generation and
all meets, are

\[
 (a,b,c)=((5\ 6),(2\ 3)(4\ 5),(1\ 2)(3\ 4)),\qquad
 ((5\ 6),(2\ 3)(4\ 5),(1\ 2)(3\ 4)(5\ 6)).
 \tag{CPI11}
\]

Reversing the first gives the source's convention \{6,5\}. The complete
degree-six certificate also verifies that the five representatives
printed in source equation (2) are pairwise inequivalent and exhaust the
five endpoint orbits:

\[
 \begin{array}{c|c}
 a&c\\\hline
 (1\ 2)&(3\ 4)\\
 (1\ 2)&(3\ 4)(5\ 6)\\
 (1\ 2)&(1\ 2)(3\ 4)\\
 (1\ 2)&(1\ 2)(3\ 4)(5\ 6)\\
 (1\ 2)(3\ 4)&(1\ 3)(2\ 4)
 \end{array}
 \tag{CPI12}
\]

Only the second and fourth endpoint orbits admit a generating middle.
The last pair represents the unique Aut(\(A_6\)) orbit. The nine
\(S_6\)-conjugacy orbits of \(S_6\) pairs fuse to five under
Aut(\(S_6\)); the two \(S_6\)-conjugacy orbits of \(A_6\) pairs fuse to
one under Aut(\(A_6\)). These observations and the table complete the
exceptional clauses of Theorems 1.1 and 1.2 and source Lemma 4.3. They
prove Theorem 3.6 and Theorem 4.4: exactly the two stated \(S_6\)
polyhedra and no \(A_6\) regular polyhedron exist.

The argument is exhaustive by UD.ExtensionOrbits: every triple moves to
a chosen endpoint orbit and then to a chosen middle orbit under its
stabilizer. Generation and every IP meet are invariant under these
moves. Each retained marking conversely is a string C-group, and source
equivalence allows exactly its automorphism and reversal orbit. The
supplied program is therefore a computational proof of the exceptional
classification, not merely a check of two constructions. It retains
GAP's exact group, automorphism and orbit algorithms; no SmallGroups or
polytope catalogue is imported. The change from the source's
computer-free exceptional analysis to this exact finite computation is
part of the proof cost.

\section{The source statements and counting
conventions}\label{the-source-statements-and-counting-conventions}

Table 1 is recovered by (CPI7), (CPI9) and the degree-six correction.
Its fourteen printed rows, including the row for degrees 1,2,3, are:

{
\begin{longtable}[]{@{}lrr@{}}
\toprule\noalign{}
Degree N & Symmetric & Alternating \\
\midrule\noalign{}
\endhead
\bottomrule\noalign{}
\endlastfoot
1,2,3 & 0 & 0 \\
4,5 & 3 & 1 \\
6 & 5 & 1 \\
7 & 9 & 2 \\
8,9 & 21 & 7 \\
10,11 & 39 & 10 \\
12,13 & 67 & 21 \\
14,15 & 105 & 28 \\
16,17 & 158 & 48 \\
18,19 & 226 & 61 \\
20 & 315 & 93 \\
30 & 1169 & 315 \\
40 & 3105 & 855 \\
50 & 6774 & 1795 \\
\end{longtable}
}

The ordinary formulas give zero below degree four, where no faithful
V-action fits. The exact formula diagnostic verifies the printed rows,
the reversible coordinates, and rational-series agreement with the floor
formulas through half-degree 500. These controls do not prove the
all-degree formulas: their proof is UD.Orbits, the preceding sign
calculation and (CPI3)--(CPI9). The whole control program is
nevertheless charged.

All thirteen numbered source items are now located: Theorems 1.1 and 1.2
in Sections 2--3; Proposition 3.1 in Section 1; Lemmas 3.2 and 3.3 in
Section 2; Definition 3.5 is the existing IP/string convention; Theorem
3.6 in Section 3; Proposition 4.1 in Section 1; Lemma 4.2 in Section 2;
Lemma 4.3 and Theorem 4.4 in Section 3. Remark 3.7 is historical
attribution to the earlier Leemans--Vauthier computation, not an
additional new mathematical target. Remark 3.4 has the exact
endpoint-reversal meaning proved above: equal-cycle endpoints form a
reversal-fixed orbit. Endpoints alone neither specify a polyhedron nor
prove a completed triple self-dual; a middle generator must also be
preserved by the relevant automorphism. We do not promote that remark to
an unconditional self-duality theorem. The printed Theorem 4.4 is read
in its section's regular-polyhedron context; the proof does not exclude
arbitrary nonregular polyhedra with group \(A_6\).

\section*{Programs and retained inputs}
The companion appendix prints the complete degree-six GAP census and the
formula diagnostic. The former is required for the exceptional classification
proved here. The latter checks the fourteen displayed rows and several
equivalent expressions; the formulas for arbitrary degree follow from the
written proof. The complete diagnostic is conservatively retained in the
program boundary. The replay driver verifies exact program bytes and runs
in a fresh directory. GAP's exact group, automorphism and orbit algorithms
are retained computational inputs.

Centralizers of permutations and finite-action classification are classical.
The ordinary ambient automorphism theorem and the string-C-group/regular-
polytope correspondence are retained. The historical earlier computation
mentioned in source Remark 3.7 is due to Leemans and Vauthier; our finite
reconstruction does not claim priority for its conclusion. The full common
paper [UD] is a dependency, not an uncharged abbreviation.

[UD] \emph{Stabilizers, extensions and the enumeration of finite geometries},
common paper in this topic. The string-C-group/polytope correspondence is
used in the standard sense of P.~McMullen and E.~Schulte,
\emph{Abstract Regular Polytopes}, Cambridge University Press, 2002.
\end{document}
